\section{Translation theorem T}
\label{sec:translation_theorem}

The next theorem translates the vanishing of the anti-invariant ledger
into the critical-line statement.  It is a positivity argument.  We
state it for the nontrivial-zero ledger of
\Cref{def:rh:nontrivial-zero-ledger}; the identical argument applies to
the ledger of represented zeros of any shell, which is how it is used
in \Cref{sec:landing_chain}.

\begin{theorem}[Translation theorem T]
\label{thm:rh:translation-T}
For any positive weights \(m_\rho\),
\[
  \AZ=0
  \quad\Longleftrightarrow\quad
  \Real\rho=\frac12\ \text{ for every }\rho\in\Znt .
\]
Consequently \(\AZ=0\) if and only if the Riemann hypothesis holds.
\end{theorem}

\begin{proof}
The two directions are
\Cref{thm:rh:translation-T-forward,thm:rh:translation-T-reverse}.  The
nontrivial zeros are exactly the zeros of \(\zeta\) in the strip
\(0<\Real s<1\), because \(\pi^{-s/2}\Gamma(s/2)\) has neither zeros
nor poles there; and every zero of \(\zeta\) other than the trivial zeros \(-2,-4,\ldots\)
lies in that strip.  For the latter, the classical facts used are
\citep{TitchmarshHeathBrown1986}: \(\zeta\) has no zeros with
\(\Real s\geq1\); \(\zeta(0)=-\frac12\); \(\Lzeta(s)=\Lzeta(1-s)\); and
\(\pi^{-s/2}\Gamma(s/2)\) never vanishes and is finite except at
\(s=0,-2,-4,\ldots\).  If \(\zeta(s)=0\) with \(\Real s\leq0\), then
\(s\neq0\), and if \(s\) is not a trivial zero the factor at \(s\) is
finite and nonzero, so \(\Lzeta(s)=0\), hence \(\Lzeta(1-s)=0\); since
\(\Real(1-s)\geq1\) and \(1-s\neq1\), the factor at \(1-s\) is finite and
nonzero, so \(\zeta(1-s)=0\), a contradiction.  Hence the critical-line
statement on \(\Znt\) is the Riemann hypothesis.
\end{proof}

\begin{theorem}[Translation theorem T, forward direction]
\label{thm:rh:translation-T-forward}
If \(\AZ=0\), then \(\Real\rho=\frac12\) for every \(\rho\in\Znt\).
\end{theorem}

\begin{proof}
Each summand \(m_\rho(\Real\rho-\frac12)^2\) is nonnegative and is at
most the extended sum \(\AZ\).  If \(\AZ=0\), each summand is therefore
zero, and since \(m_\rho>0\) we get \((\Real\rho-\frac12)^2=0\).
\end{proof}

\begin{theorem}[Translation theorem T, reverse direction]
\label{thm:rh:translation-T-reverse}
If \(\Real\rho=\frac12\) for every \(\rho\in\Znt\), then \(\AZ=0\).
\end{theorem}

\begin{proof}
Every summand vanishes, so every partial sum vanishes, and so does
their supremum.
\end{proof}

The equivalence between \(\AZ=0\) and the critical-line statement uses
nothing about the zeros beyond positivity of the weights; only its
identification with the Riemann hypothesis uses the analytic facts
recalled in the proof.  Its role is to let the rest of the paper work with the single
quantity \(\AZ\), or with its finite window sums, instead of with the
zero set.
